%%%%%%%%%%%%%%%%%%%%%%%%%%%%%%%%%%%%%%%%%%%%%%%%%%%%%%%
%%%%%%%%%%%%%%%%%%%%%%%%%%%%%%%%%%%%%%%%%%%%%%%%%%%%%%%
%%
%%       Template de TCC (Monografia) da Universidade Estadual da Paraíba
%%					Centro de Ciências e Tecnologia
%%					Departamento de Matemática
%%					Curso de Licenciatura em Matemática
%%          Contato: coord.matematica.cct@setor.uepb.edu.br
%%
%%%%%%%%%%%%%%%%%%%%%%%%%%%%%%%%%%%%%%%%%%%%%%%%%%%%%%%
%%%%%%%%%%%%%%%%%%%%%%%%%%%%%%%%%%%%%%%%%%%%%%%%%%%%%%%

\documentclass[12pt, a4paper]{report}
\usepackage[top=3cm, bottom=2cm, right=2cm, left=3cm]{geometry}
\usepackage[utf8]{inputenc}
\usepackage[greek.polutoniko, brazil]{babel}
\usepackage[final]{pdfpages}
\usepackage{csquotes}

\usepackage[labelsep=endash]{caption}
\renewcommand{\thefigure}{\arabic{figure}}
\usepackage{indentfirst}
\usepackage{blindtext} 
\usepackage[titletoc]{appendix}
\usepackage[round]{natbib}  % citacao tipo texto e tipo parentes
\addto{\captionsbrazil}{\renewcommand{\bibname}{\normalsize \bf \hfill REFERÊNCIAS \hfill}}
\pagestyle{myheadings}
\usepackage{epigraph}
\renewcommand{\textflush}{flushepinormal}
\usepackage{graphicx}
\usepackage{float}
\usepackage{tikz}
\usepackage[all]{xy}
\usepackage{capt-of}
\usetikzlibrary{arrows.meta,bending}
\usepackage{amsmath, amssymb, amsfonts, amsthm}
\usepackage{MnSymbol}
\usepackage{thmtools}
\DeclareMathOperator{\mdc}{mdc}
\DeclareMathOperator{\sen}{sen}
\usepackage{physics}
\usepackage{hyphenat}
\usepackage{cancel}
\usepackage{setspace}
\linespread{1.25}
\usepackage{boxedminipage}
\usepackage{enumitem}

%-----------------------------------
\usepackage{nomencl}
\makenomenclature
\setlength{\nomlabelwidth}{3cm}
\setlength{\nomitemsep}{-\parsep}

\renewcommand{\nomname}{\normalsize \makebox[\linewidth]{\MakeUppercase{Lista de símbolos}}}
%--------------------------------------------------
\usepackage{datetime}


\usepackage[pdftex]{hyperref}

% partes do documento

\usepackage{titlesec}
\titleformat*{\section}{\normalsize \bfseries}
\titleformat*{\subsection}{\normalsize}
\titleformat*{\subsubsection}{\normalsize \bfseries}
\titleformat*{\paragraph}{\large\bfseries}
\titleformat*{\subparagraph}{\large\bfseries}
\titleformat{\chapter}{\bfseries}{\filright \thechapter}{20pt}{\filright}
\titlespacing*{\chapter}{0pt}{0pt}{*1.5}
\titlespacing{\section}{0pt}{*1.5}{*1.5}

% formatação do sumário de acordo com a abnt nbr 6027: 2012
% -------------------------------------------------------
\usepackage{tocloft}
\renewcommand\cftchapnumwidth{3.8em}
\renewcommand\cftsecnumwidth{3.8em}
\renewcommand\cftsubsecnumwidth{3.8em}
\renewcommand{\cftsecfont}{\bfseries }
%\renewcommand{\cftsubsecfont}{}
\setlength{\cftsecindent}{0pt}% Remove indent for \section
\setlength{\cftsubsecindent}{0pt}%

%\renewcommand{\cftdot}{\ensuremath{\ast}}
\renewcommand{\cfttoctitlefont}{\hfil \normalfont \centering \bfseries\MakeUppercase}



\renewcommand{\cftaftertoctitle}{
\\[\baselineskip]\mbox{}\hfill{\normalfont Página}}

\setlength{\cftbeforetoctitleskip}{-1.5em}
\setlength{\cftaftertoctitleskip}{0pt}

%-----------------------------------------

\newtheorem{theorem}{Teorema}[chapter]
\newtheorem{teorema}{Teorema}[chapter]
\newtheorem{corollary}{Corolário}[theorem]
\newtheorem{lemma}[theorem]{Lema}

\theoremstyle{remark}
\newtheorem{remark}{Observação}[chapter]

\theoremstyle{definition}
\newtheorem{definition}{Definição}[chapter]

\theoremstyle{definition}
\newtheorem{exemplo}{Exemplo}[chapter]

\theoremstyle{plain}
\newtheorem{lema}{Lema}[chapter]

\theoremstyle{plain}
\newtheorem{prop}{Proposição}[chapter]


\theoremstyle{plain}
\newtheorem{cor}{Corolário}[chapter]

\declaretheorem[style=definition,qed=$\filledmedtriangleleft$, name=Aplicação, numberwithin=chapter]{ap}

\theoremstyle{definition}
\newtheorem*{nota}{Notação}
%-----------------------------------

\begin{document}
%-------
%---CAPA

\begin{titlepage}
\begin{center}
\begin{figure}
    \centering
    \includegraphics[height=!, width= 2cm]{imagens/uepb.png}
\end{figure}

\textbf{UNIVERSIDADE ESTADUAL DA PARAÍBA\\
CAMPUS X\\
CENTRO XXXXXXXXX\\
DEPARTAMENTO DE XXXXXXXX\\
CURSO DE XXXXXXXXX}

\vspace{4cm}


\textbf{NOME DO ESTUDANTE}   

\vspace{4cm}
\textbf{TÍTULO DO TRABALHO}

\vfill

\textbf{CIDADE\\
ANO} 
    
\end{center}
\end{titlepage}

%----------------
%--FOLHA DE ROSTO

\thispagestyle{empty}
\begin{center}
 
    
\textbf{NOME DO ESTUDANTE}    

\vspace{5cm}

\textbf{TÍTULO DO TRABALHO}   

\vspace{5cm}

\begin{flushright}
\parbox{8cm}{
\singlespacing{\begin{hyphenrules}{nohyphenation}
Trabalho de Conclusão de Curso apresentado ao Departamento de XXXXXXX do Centro de XXXXXXXXXX da Universidade Estadual da Paraíba como requisito parcial à obtenção do título de Licenciado(a) em XXXXXXX.        % natureza do trabalho
\vspace{0.5cm}    

\textbf{Área de concentração}:   xxxxxxxxxxxxx
\end{hyphenrules}}}
\end{flushright}


\vspace{2cm}

\textbf{Orientador}: Prof. xxxxxxxxxxx
\vfill


\textbf{CIDADE\\
ANO}
\end{center}
\newpage

%----------------
%--FOLHA DE ROSTO

\thispagestyle{empty}
\begin{center}
\textbf{NOME DO ALUNO}

\vspace{1cm}

\textbf{TÍTULO DO TRABALHO}

\end{center}

\vspace{.4cm}
\begin{flushright}
\parbox{8cm}{
\singlespacing{\begin{hyphenrules}{nohyphenation}
Trabalho de Conclusão de Curso apresentado ao Departamento de XXXXXXXX do Centro de XXXXXXXX da Universidade Estadual da Paraíba como requisito parcial à obtenção do título de Licenciado(a) em XXXXXXX.
\vspace{0.5cm}

\textbf{Área de concentração}:  xxxxxxxxxx
\end{hyphenrules}}}
\end{flushright}

\vspace{.6cm}

\noindent Aprovado em: ---------

\vspace{1.5cm}

\begin{center}
\textbf{BANCA EXAMINADORA}
\vspace{1.5cm}

\noindent\rule{12cm}{0.4pt} \\
Prof. Dr. xxxxxxxxxxxxxx (Orientador)\\
Universidade Estadual da Paraíba (UEPB)
\vspace{2cm}

\noindent\rule{12cm}{0.4pt}\\
Profa. Dra. xxxxxxxxxxxxxx  \\
Universidade Estadual da Paraíba (UEPB)
\vspace{2cm}

\noindent\rule{12cm}{0.4pt}\\
Prof. Dr. xxxxxxxxxxxxxxx \\
Universidade XXXXXX (CASO MEMBRO EXTERNO)
\vspace{2cm}

\end{center}

%--------------
%---DEDICATÓRIA

\thispagestyle{empty}
\null
\vfill
\begin{flushright}
\parbox{4cm}{
\singlespacing{\begin{hyphenrules}{nohyphenation}
Dedicatória.
\end{hyphenrules}
}}
\end{flushright}
\thispagestyle{empty}
\newpage

%-----------------
%---AGRADECIMENTOS

\thispagestyle{empty}
\section*{\normalsize \centering \textbf{AGRADECIMENTOS}}

Escreva aqui os seus agradecimentos...

\newpage

%-----------
%---EPÍGRAFE

\thispagestyle{empty}
\null
\vfill
\begin{flushright}
\parbox{8cm}{
\singlespacing{\begin{hyphenrules}{nohyphenation}
``Epígrafe - Citação relacionada com o tema do trabalho, com indicação de autoria. Dever seguir as diretrizes da NBR 10.520/2002 de Citação em documentos.''
\end{hyphenrules}
}}
\end{flushright}
\newpage

%---------
%---RESUMO

\thispagestyle{empty}
\section*{\normalsize \centering \textbf{RESUMO}}

\noindent  O resumo deverá ser uma apresentação concisa dos pontos relevantes de um documento. Deve ressaltar o objetivo, o método, os resultados e as conclusões do documento. Deverá ser composto de uma sequência de frases concisas, afirmativas e não de enumeração de tópicos. Recomenda-se o uso de parágrafo único. Deve-se usar o verbo na voz ativa e na terceira pessoa do singular.

\vspace*{0.5cm}

\noindent\textbf{Palavras-chave:} Palavras-chave 1. Palavras-chave 2. Palavras-chave 3.
 % no máximo 4 palavras-chave separadas entre si por ponto
 \newpage

%---------
%---ABSTRACT

\thispagestyle{empty}
\section*{\normalsize \centering \textbf{ABSTRACT}}

\noindent Resumo em inglês

\vspace*{0.5cm}

\noindent\textbf{Keywords:} Keyword 1. Keyword 2. Keyword 3.
\newpage

%----------
%---SUMÁRIO

\tableofcontents \thispagestyle{empty} 

%----------
%---INTRODUÇÃO

\chapter{INTRODUÇÃO}
\thispagestyle{myheadings}

Texto, texto texto...

Texto, texto texto...

%----------
%---CAPÍTULO I

\chapter{TEORIA DOS NÚMEROS}
\thispagestyle{myheadings}

Texto, texto texto...

Texto, texto texto...

\section{Seção primária}

Texto, texto texto...

Texto, texto texto...

\subsection{Seção secundária}

Texto, texto texto...

Texto, texto texto...

%----------
%---CAPÍTULO II

\chapter{TÍTULO DO CAPÍTULO}
\thispagestyle{myheadings}

Texto, texto texto...

Texto, texto texto...

\section{Isto é uma seção}

Texto, texto texto...

Texto, texto texto...

\subsection{Isto é uma subseção}

Texto, texto texto...

Texto, texto texto...

%----------
%---CAPÍTULO III

\chapter{TÍTULO DO CAPÍTULO}
\thispagestyle{myheadings}

Texto, texto texto...

Texto, texto texto...

\section{Isto é uma seção}

Texto, texto texto...

Texto, texto texto...

\subsection{Isto é uma subseção}

Texto, texto texto...

Texto, texto texto...

%--------------
%---CAPÍTULO VI

\chapter{CONCLUSÃO}
\thispagestyle{myheadings}

Texto, texto texto...

Texto, texto texto...

\section{Isto é uma seção}

Texto, texto texto...

Texto, texto texto...

\subsection{Isto é uma subseção}

Texto, texto texto...

Texto, texto texto...

%-----------------------------
%---REFERÊNCIAS BIBLIOGRÁFICAS

\addcontentsline{toc}{chapter}{\protect\numberline{}REFERÊNCIAS}
\begin{thebibliography}{20}
\thispagestyle{myheadings}

\bibitem[Alencar Filho(1981)]{alencarfilho1981}
ALENCAR FILHO, Edgard de. \textbf{Teoria elementar dos números}. São Paulo: Nobel, 1981.
exemplo de citação: (ALENCAR FILHO, 1981)
exemplo de citação: Alencar Filho, 1981

Inserir exemplo de referência de Artigo





% exemplo de referência

\end{thebibliography}

%------------
%---APÊNDICES

\addcontentsline{toc}{chapter}{\protect\numberline{}APÊNDICE A -- TÍTULO DO APÊNDICE}
\addtocontents{toc}{\protect\setcounter{tocdepth}{0}}
\appendix
\setcounter{chapter}{0}
\renewcommand{\thechapter}{\Alph{chapter}}
\chapter*{\hfill APÊNDICE A -- TÍTULO DO APÊNDICE \hfill}
\stepcounter{chapter}
\thispagestyle{myheadings}

Conteúdo do Apêndice A...

Conteúdo do Apêndice A...


\end{document}
